% Overleaf-Datei 01_Introduction_V2.tex, Spiegel in docs/overleaf/.
% Stand: Overleaf-Fassung vom 2026-09-14 plus die dort noetigen Korrekturen;
% Schluss ab "overfitting" aus docs/background.tex uebernommen.
% Nur zusammen mit der Overleaf-Datei aendern.

% Introduction der Thesis, Entwurf 2026-09-12. Zwei Teile: (1) die Motivation
% des Delta-Kapitels, aus chapter_delta.tex \section{Motivation} hierher
% gezogen (Absaetze 1-4 dort; Gleichungen und Architektur bleiben im
% Kapitel); (2) das Setting des OMol25-Kapitels in der zuletzt vereinbarten
% Fassung. Labels: ch:background, ch:delta, ch:mr, sec:mace-training,
% sec:delta-architecture, sec:delta-data, sec:mr-methods, app:marks.

\chapter{Introduction}
\label{ch:introduction}

% ---------------------------------------------------------------- Teil 1
% Wortlaut aus chapter_delta.tex \section{Motivation}, Absaetze 1-4.
% Geaendert nur: Platzhalter aufgeloest, Praemisse benannt, Gleichungen und
% Architektur im Kapitel belassen, Roadmap angehaengt.

Machine-learned interatomic potentials (MLIPs) such as PaiNN or MACE
are commonly trained on DFT calculations. Trained on reaction paths,
they can search transition states and predict barriers at a fraction
of the cost of DFT. Transition1x~\cite{schreiner2022} was the first
large dataset built for this purpose: 9.6 million DFT energies and
forces on and around the reaction paths of 10\,073 organic reactions,
where earlier datasets held equilibrium structures and their
surroundings. It demonstrated that a PaiNN model trained on these
${\sim}10$ million reaction-path geometries reproduces DFT energies
and forces. However,
Transition1x was computed at \mbox{$\omega$B97X/6-31G(d)}, a
comparatively cheap level of theory chosen to make the generation of
ten million data points computationally feasible. More accurate
levels of theory exist. The range-separated
meta-GGA $\omega$B97M-V with a triple-$\zeta$ basis is among the most
accurate DFT methods for barrier heights~\cite{mardirossian2016,
goerigk2017}. Any model trained on the Transition1x dataset can only
be as accurate as the level of theory the dataset was generated with,
so upgrading the level of theory is desirable.

Such an upgrade does not necessarily require repeating the NEB searches at the
higher level of theory. A dataset like Transition1x consists of two
things: the geometries, which sample the configuration space along
reaction paths, and the energy and force labels attached to them. The
premise is that the geometries remain valid at any level of theory:
an MLIP only needs them to cover the relevant region, which the
Transition1x paths already do. Only the labels are tied to the cheap
level. Replacing them requires one single-point calculation per geometry at
the higher level; this is what relabelling means. It is far cheaper
than relaxing the paths again at that level, which would take hundreds
of gradient calculations per reaction instead of one per geometry
kept.

OMol25~\cite{omol25} relabelled the entire Transition1x dataset at
\mbox{$\omega$B97M-V/def2-TZVPD}\footnote{def2-TZVPD adds diffuse
functions to def2-TZVP; this work omits them for cost.} using
unrestricted Kohn--Sham\footnote{Chapter~\ref{ch:delta} uses restricted
Kohn--Sham throughout. The implications of unrestricted versus
restricted Kohn--Sham are discussed in Chapter~\ref{ch:mr}.}.

In Chapter~\ref{ch:delta}, we pursue a different goal. Relabelling a
whole dataset costs one calculation per geometry at the higher level,
and any way to reach that level at a fraction of the cost is
desirable. We therefore ask whether it can be reached by relabelling
only a small part of the data: about 80\,000 geometries, under 1\,\%
of the dataset. A MACE model is trained on Transition1x to predict the cheap level,
\mbox{$\omega$B97X/6-31G(d)} (Section~\ref{sec:mace-training}), and
a small correction head on top of it
(Section~\ref{sec:delta-architecture}) learns the difference between
the cheap and the expensive level, \mbox{$\omega$B97M-V/def2-TZVP}. This is feasible with so little data because
the difference is nearly constant within a reaction
(Section~\ref{sec:delta-data}): a smooth offset is far easier to learn
than the potential energy surface itself. Base model and correction
head combine into one corrected potential, MACE+$\Delta$, which can be
used wherever MACE can: for single-point predictions or as the
potential driving a NEB optimisation. The chapter evaluates it on 30
test reactions grouped by multireference (MR) character, at fixed
geometries and in searches the model drives itself. It finds that the force errors drop to a third on all reactions, and
that the correction reproduces the higher level where the
multireference character is low: there the barrier errors drop to a
third as well, while on the mid- and high-MR reactions the correction
overshoots.

% ---------------------------------------------------------------- Teil 2

Chapter~\ref{ch:delta} raises the level of theory of the labels and
leaves the geometries unchanged, on the premise that changing the labels is
enough to lift the model to the higher level of theory. That is how the field raises fidelity: delta
learning~\cite{ramakrishnan2015}, multi-fidelity
training~\cite{messerly2025}, dataset merging~\cite{shiota2024} and
reactive potentials with semi-empirical geometries~\cite{ren2026} all
supply labels at a higher level of theory on geometries they do not
re-optimise. OMol25~\cite{omol25} went one step further: it took
Transition1x, a dataset generated restricted at a cheap level of
theory, and relabelled the full set, not a fraction of it,
unrestricted at an expensive one, again without re-optimising a single
path. Two questions lie underneath this
practice, and this thesis asks both:

\noindent
\begin{minipage}{\linewidth}
\begin{description}
  \item[Research question 1 (RQ1)] \textbf{Can a model be lifted to a higher
    level of theory by relabelling only a subset of its training data?}
  \item[Research question 2 (RQ2)] \textbf{Can a model learn a surface from
    paths that were never relaxed on that surface?}
\end{description}
\end{minipage}

Chapter~\ref{ch:delta} answers RQ1. For that chapter, RQ2 needs no
test: both levels of theory are restricted, and when the transition
states are searched with DFT at either level, the two results lie
within 0.007\,\AA{} of each other. A path relaxed at the cheap level
therefore already covers the immediate region the transition state
lies in when calculated with the expensive method, and relabelling the
path without re-optimising it is enough.

This stops being true when the spin formalism changes as well. At a
broken-symmetry transition state the restricted solution is unstable
(Section~\ref{sec:background-stability}), and the transition state
found unrestricted can lie far from the one found restricted. A path
relaxed restricted then no longer covers the region the transition
state lies in when calculated unrestricted, and RQ2 becomes a real
question. OMol25 is exactly this case: every label recomputed
unrestricted at \mbox{$\omega$B97M-V/def2-TZVPD}, on the inherited
Transition1x geometries and on many structures beyond
them~\cite{omol25}. Chapter~\ref{ch:mr} asks RQ2 there.

RQ1 was the task this thesis was set. RQ2 arose during that work,
when the OMol25 models, used along the way as a point of comparison,
delivered structures other than the restricted reference transition
states this work had built and used for Chapter~\ref{ch:delta}, on
some reactions. Whether these models can find broken-symmetry
transition states at all had been tested nowhere, because the
published benchmarks only use restricted references, as the rest of
this chapter shows; the Preface gives the account.

The published benchmarks that test the OMol25 models at locating
transition states report high success rates. According to them, four
OMol25 models, three
of them used in Chapter~\ref{ch:mr}, reach the reference saddle in 91.8\,\% of
automated transition-state searches, and an MLIP search is recommended
as the default first stage of such searches~\cite{marks2026}. Hait et
al. use eSEN to produce the starting guess for a DFT transition-state
optimisation, and the DFT optimisation converges from that guess in
fewer steps than from a guess produced with DFT alone~\cite{hait2025}.

What these benchmarks do not measure is the outcome at broken-symmetry
transition states, because both test against transition states
optimised on the restricted surface. Hait et al. say so: they
re-optimised their 121 references unrestricted, 27 broke spin symmetry
and moved away from the supplied saddle, and all 27 were excluded as
defective references rather than replaced~\cite{hait2025}. The authors
expected models like eSEN to be challenged by these cases even with
unrestricted training data, and did not test it. Marks et al. report
nothing about the SCF solution of their references, so we ran the
stability analysis ourselves, at the paper's own level of theory. Of
the 25 even-electron references, 24 are closed-shell and one is
broken-symmetry (Appendix~\ref{app:marks}). That one was still optimised restricted, so it is not a transition
state of the unrestricted surface and should not serve as a
reference. In the work of Marks et al.\ it still serves as one, and so
a model that finds the same restricted structure is scored as a
success, although the true transition state lies elsewhere.

The structures are not rare. OMol25 reports $\langle S^2 \rangle >
0.001$, that is, broken-symmetry solutions, for just under 5\,\% of
its unrestricted Transition1x configurations~\cite{omol25}. That number counts every configuration of
a reaction path, most of which lie near reactant and product where the
restricted solution is stable, so the share at the saddle points
themselves is not known and can be assumed to be higher. In the
reference set of Hait et al. it is a fifth~\cite{hait2025}.

Where a transition state is an open-shell singlet, quantum chemistry
has a rule. The affordable treatment is broken-symmetry unrestricted
DFT~\cite{grafenstein2002}, and the geometry is optimised on that
surface, because the restricted saddle is not a stationary point of
it~\cite{crawford2001}. The label should be computed on the surface the
geometry was optimised on. The
benchmarks above and the training data behind the OMol25 models each
depart from this rule, in different ways (Table~\ref{tab:surfaces}).

\begin{table}[h]
\centering
\small
\caption{On which surface the geometry and the label of a
broken-symmetry transition state lie.}
\label{tab:surfaces}
\begin{tabular}{@{}lll@{}}
\toprule
 & geometry & label \\
\midrule
rule & broken-symmetry & broken-symmetry \\
benchmark references (Appendix~\ref{app:marks}, \cite{hait2025, marks2026}) & restricted & restricted \\
OMol25 training data~\cite{omol25} & restricted & broken-symmetry \\
\bottomrule
\end{tabular}
\end{table}

The benchmark references are restricted throughout, geometry and label,
which is why a model that reproduces the restricted saddle matches them
(Table~\ref{tab:surfaces}).
The OMol25 training data keeps half of the rule: the Transition1x
geometries come from restricted searches that never tested the
solution~\cite{schreiner2022}, and OMol25 recomputed their labels
unrestricted and left the geometries where they were~\cite{omol25}.
Whether half is enough has not been measured.

Chapter~\ref{ch:mr} attempts to close that gap. From the Transition1x
test split it selects reactions of high multireference character,
where broken-symmetry solutions are expected at the transition state,
together with a mid- and a low-MR group as control. The OMol25 models
UMA-S, UMA-M and eSEN then drive the NEB searches, and every transition
state they find is re-evaluated with unrestricted DFT, the protocol
OMol25 used for its labels, extended by a stability analysis. The
stability analysis decides whether the lowest solution at the structure
the model found is closed-shell or broken-symmetry, and it splits the
transition states into these two groups. On the two groups the same
metrics are applied, to test whether the models are equally successful
at finding a transition state in a broken-symmetry region as in a
closed-shell one.

They are not. At broken-symmetry transition states the models deliver
structures that are not stationary points of the unrestricted surface,
and their forces there deviate more from DFT than at closed-shell ones.
The chapter traces this to the training data. The training geometries of the broken-symmetry reactions are not
stationary points of the unrestricted surface: the force there is
not zero, so the models have never seen a transition state of that
surface, only the region around it.

\newpage
\chapter{Background and Theory}

\section{Potential Energy Surfaces}
\label{sec:background-pes}

\subsection{The Born--Oppenheimer Surface}
\label{sec:background-bo}

Nuclei are thousands of times heavier than electrons. On the timescale of
nuclear motion the electrons adjust instantly, so the electronic problem can
be solved with the nuclei held fixed at positions $\mathbf{R}$,
\begin{equation}
  \hat{H}_\text{el}(\mathbf{R})\, \psi_k(\mathbf{r};\mathbf{R})
  = E_k(\mathbf{R})\, \psi_k(\mathbf{r};\mathbf{R}),
  \qquad E_0 \le E_1 \le \dots
  \label{eq:electronic-schrodinger}
\end{equation}
The lowest eigenvalue $E_0(\mathbf{R})$, computed for every nuclear
arrangement, is the potential energy surface (PES). The nuclei move on it and
feel the force $-\nabla E_0(\mathbf{R})$. The approximation holds when the
ground state is well separated from the excited states.

\subsection{Stationary Points}
\label{sec:background-stationary}

The chemistry of a surface lies in its stationary points, where the gradient
$\mathbf{g} = \nabla_\mathbf{R} E$ vanishes. The eigenvalues of the Hessian
$\mathbf{H}_{ij} = \partial^2 E / \partial R_i \partial R_j$ tell them apart:

\begin{itemize}
  \item \textbf{Local minima} ($\mathbf{H}$ positive definite) are stable or
    metastable structures --- the reactants and products of a reaction.
  \item \textbf{First-order saddle points} (exactly one negative eigenvalue)
    are \textbf{transition states} (TS): the highest point on the
    lowest-energy path between two minima. The eigenvector of the negative
    eigenvalue is the reaction coordinate at the TS.
\end{itemize}

The energy difference between transition state and reactant is the barrier
$\Delta E^\ddagger$. Transition state theory relates it to the rate constant
by $k \propto \exp(-\Delta E^\ddagger / k_\text{B} T)$~\cite{eyring1935}, so
an error in the barrier enters the rate exponentially. An error of
1~kcal/mol (43~meV) changes the rate at room temperature by a factor of
five; this value is called chemical accuracy.

\subsection{Geometry and Energy}
\label{sec:geom-energy}

A barrier calculation has two parts: the geometry of the transition state,
and the energy of the surface at that geometry relative to the reactant. The
two can be computed with different methods. A geometry optimised with one
method can be used for an energy calculation with another, a single-point
calculation. A reference calculation therefore has a geometry level and an
energy level. Common practice lets the geometry level be the lower one,
because structures depend less on functional and basis set than
energies~\cite{bursch2022}.

\section{Electronic Structure Methods}
\label{sec:background-electronic}

Eq.~\eqref{eq:electronic-schrodinger} can be solved exactly for one
electron only. To make it solvable for a molecule, every method
approximates the wavefunction by an ansatz, an assumed form fixed
before the calculation starts. The ansatz says what the electrons are
allowed to do: for example, that each electron occupies one orbital
and feels the others only as an average, or that a spin-up and a
spin-down electron share the same orbital. The calculation then finds
the lowest energy the ansatz allows~\cite{szabo1996}. If the true
ground state does something the ansatz forbids, the calculation still
converges and reports an energy and a gradient, and nothing in the
output says that a lower state exists. A calculation can only tell how
well it did within its ansatz, not whether the ansatz is a
sufficiently accurate description of the molecule. That takes a
separate test (Section~\ref{sec:background-stability}).

% =====================================================================
\subsection{Single-Reference and Multireference Systems}
\label{sec:sr-vs-mr}

Electrons occupy orbitals, at most two per orbital with opposite spin. One
assignment of all electrons to orbitals is a configuration. As a
wavefunction, a configuration is a Slater determinant $\Phi_I$. The exact
electronic state is a weighted sum of all configurations,
%
\begin{equation}
  \Psi = \sum_I c_I \, \Phi_I , \qquad \sum_I |c_I|^2 = 1 .
  \label{eq:ci-expansion}
\end{equation}
%
Two cases are distinguished by how the weights $|c_I|^2$ are distributed:

\begin{itemize}
  \item \textbf{Single-reference.} One configuration has a weight close to
    one; all others are small corrections. This is the case for most
    molecules near their equilibrium geometry, where all electrons are paired
    in bonding orbitals.
  \item \textbf{Multireference.} Two or more configurations have comparable
    weight. Example: a half-broken bond. The two electrons of the bond can
    be described as a pair in the bonding orbital, or as one electron on
    each fragment. Both configurations have similar energy, and the state is
    a mixture of both. This occurs at breaking bonds, at diradicals, and at
    many transition states.
\end{itemize}

Methods are built for one of the two cases:

\begin{itemize}
  \item \textbf{Single-reference methods} start from one determinant:
    Hartree--Fock, Kohn--Sham DFT, coupled cluster.
  \item \textbf{Multireference methods} optimise several determinants at
    once: CASSCF, NEVPT2 (Section~\ref{sec:background-beyond}).
\end{itemize}

The following sections treat single-reference methods.


\subsection{Kohn--Sham Density Functional Theory}
\label{sec:background-dft}

Density functional theory (DFT) replaces the wavefunction, a function of all
$3N_e$ electron coordinates, by the electron density $\rho(\mathbf{r})$, a
function of three coordinates. The ground-state energy is a functional of
the density alone~\cite{hohenberg1964}. This makes DFT the most widely used
electronic structure method: it reaches useful accuracy at a cost of
$\mathcal{O}(N^3)$ to $\mathcal{O}(N^4)$ in the number of basis functions.

In the Kohn--Sham formulation~\cite{kohn1965} the density is built from one
determinant of orbitals $\phi_i$, and the energy is split into four terms:
%
\begin{equation}
  E[\rho] = T_s[\rho] + E_\text{ext}[\rho] + E_\text{H}[\rho] + E_\text{XC}[\rho] .
  \label{eq:ks-energy}
\end{equation}

\begin{itemize}
  \item $T_s$: kinetic energy of the electrons in the orbitals. Exact.
  \item $E_\text{ext}$: attraction between electrons and nuclei. Exact.
  \item $E_\text{H}$: classical repulsion of the density with itself. Exact.
  \item $E_\text{XC}$: exchange and correlation, everything the other three
    terms leave out. Its exact form is unknown; it is approximated by an
    exchange--correlation functional.
\end{itemize}

The orbitals are found by the self-consistent field (SCF) iteration:

\begin{enumerate}
  \item a starting guess for the orbitals gives a density;
  \item the density gives an effective potential;
  \item the potential gives new orbitals;
  \item repeat until the density no longer changes.
\end{enumerate}

The converged result is a stationary point of the energy with respect
to the orbitals, and there can be more than one. Which one the
iteration ends in depends on the starting guess, and nothing in the
iteration checks whether a lower one exists. A converged calculation
can therefore sit above the lowest solution without any sign of it.
The test for that is the stability analysis of
Section~\ref{sec:background-stability}.
%%%%%%%%%%%%%%%%%%%%%%%%%%%%%%%%%%%%%%%%%%%%%%
Equation~\eqref{eq:ks-energy} and the SCF iteration leave three things
unspecified (Figure~\ref{fig:choices}). The term $E_\text{XC}$ needs a formula; this is the choice of
functional (Section~\ref{sec:background-functionals}). The orbitals
$\phi_i$ need a mathematical form in which they can be stored and
optimised; this is the choice of basis set
(Section~\ref{sec:background-basis}). And each orbital carries a spin;
whether spin-up and spin-down electrons must share the same orbital or may
have different ones is the choice of spin formalism
(Section~\ref{sec:background-spin}). The first two are named in the level
of theory, written functional/basis. 

\begin{figure}[h]
\centering
\resizebox{\linewidth}{!}{%
\begin{tikzpicture}[
  box/.style={draw, rounded corners=3pt, align=center, font=\small,
              inner sep=5pt, text width=4.4cm, minimum height=1.7cm},
  term/.style={inner sep=1.5pt, anchor=west},
  tgt/.style={inner sep=3pt, rounded corners=2pt},
  arr/.style={->, thick, >=stealth, shorten >=3pt, shorten <=3pt},
]
% --- line 1: energy, unknown term E_XC
\node[term] (l1) at (0,0)
  {$E[\rho] \;=\; T_s + E_\text{ext} + E_\text{H} \;+$};
\node[tgt, fill=blue!12, base right=3pt of l1] (exc) {$E_\text{XC}[\rho]$};
\node[font=\footnotesize, anchor=west] at (0,-0.55)
  {$T_s$, $E_\text{ext}$, $E_\text{H}$ exact for the given orbitals};
% --- rows 2 and 3 at fixed distance below the centre of E_XC
\coordinate (row2) at ($(l1.west |- exc.center)+(0,-2.1)$);
\coordinate (row3) at ($(l1.west |- exc.center)+(0,-4.2)$);
% --- line 2: orbitals from basis functions
\node[term] (l2) at (row2) {$\phi_i \;=\; \sum_\mu c_{\mu i}\,$};
\node[tgt, fill=green!12, right=0pt of l2] (chi) {$\chi_\mu$};
% --- line 3: density from the orbitals of both spins
\node[term] (l3) at (row3)
  {$\rho \;=\; \sum_i \big|\phi_i^{\alpha}\big|^2 + \big|\phi_i^{\beta}\big|^2$,
   \quad};
\node[tgt, fill=orange!15, right=0pt of l3] (sp)
  {$\phi_i^{\alpha} \overset{?}{=} \phi_i^{\beta}$};
\begin{scope}[on background layer]
  \node[draw, rounded corners=3pt, fill=gray!6, inner sep=8pt,
        fit=(l1)(exc)(l2)(chi)(l3)(sp)] (eq) {};
\end{scope}

% --- the three choices, at the same fixed rows
\node[box, fill=blue!6,   anchor=west] (func)  at ($(eq.east |- exc.center)+(1.2,0)$)
  {\textbf{Functional}\\ the formula for $E_\text{XC}$\\ \emph{e.g.\ $\omega$B97M-V}};
\node[box, fill=green!6,  anchor=west] (basis) at ($(eq.east |- exc.center)+(1.2,-2.1)$)
  {\textbf{Basis set}\\ the fixed functions $\chi_\mu$ the orbitals are built from\\ \emph{e.g.\ def2-TZVP}};
\node[box, fill=orange!8, anchor=west] (spin)  at ($(eq.east |- exc.center)+(1.2,-4.2)$)
  {\textbf{Spin formalism}\\ $\alpha$, $\beta$ forced to share orbitals\\ (RKS) or not (UKS)};

\draw[arr] (func.west)  -- (exc.east);
\draw[arr] (basis.west) -- (chi.east);
\draw[arr] (spin.west)  -- (sp.east);

% --- brace: the level of theory names the first two
\draw[decorate, decoration={brace, amplitude=5pt}, thick]
  ([xshift=4pt]func.north east) -- ([xshift=4pt]basis.south east)
  node[midway, right=7pt, font=\small, align=left]
  {level of\\ theory};
\end{tikzpicture}%
}
\caption{The three choices of a Kohn--Sham calculation and where they enter.
The energy depends on the density through the unknown term $E_\text{XC}$,
which the functional supplies. The orbitals are built from the basis
functions $\chi_\mu$, which the basis set supplies. The density is built
from the orbitals of both spins, and the spin formalism decides whether
these orbitals are the same. The level of theory, written functional/basis,
names the first two choices.}
\label{fig:choices}
\end{figure}

\subsection{Exchange--Correlation Functionals}
\label{sec:background-functionals}

The presentation in this and the following section follows
Jensen~\cite{jensen2017} and Koch and Holthausen~\cite{koch2001}.

The exchange--correlation energy $E_\text{XC}$ collects the quantum
mechanical effects that the other three terms of the Kohn--Sham energy do
not describe: the exchange interaction between electrons of the same spin,
and the correlation between all electrons. Its exact dependence on the
density is unknown. An approximate expression for $E_\text{XC}$ as a
function of the density is called an exchange--correlation functional. The
functional is the only approximation in a Kohn--Sham calculation, and the
quality of the result depends on it.

Functionals are classified by the information they use
(Table~\ref{tab:functionals}). The simplest use only the density at each
point. Better ones add the gradient of the density, then the kinetic energy
density of the orbitals. The next step concerns the exchange energy, the
lowering of the electron repulsion that follows from the Pauli principle:
electrons of the same spin avoid each other. This energy can be
approximated from the density, or calculated exactly from the orbitals with
the Hartree--Fock expression. Hybrid functionals use a mixture of both, with
typically 20 to 25\,\% exact exchange. This ordering is known as Jacob's
ladder. Each step increases the computational cost and, on average, the
accuracy.

\begin{table}[h]
\centering
\caption{Jacob's ladder of exchange--correlation
functionals~\cite{perdew2001}. Each rung uses the ingredients of the rungs
below and adds one.}
\label{tab:functionals}
\begin{tabular}{@{}l p{7.4cm} l@{}}
\toprule
rung & adds & example \\
\midrule
LDA & the density $\rho$ at each point & SVWN \\
GGA & the gradient of the density, $\nabla\rho$ & PBE \\
meta-GGA & the kinetic energy density of the occupied orbitals, $\tau$,
  and/or the Laplacian of the density, $\nabla^2\rho$ & TPSS \\
hybrid & a constant fraction of exact exchange, computed from the occupied
  orbitals as in Hartree--Fock theory & B3LYP \\
range-separated hybrid & the same, but the fraction of exact exchange
  depends on the interelectronic distance: small at short range, full at
  long range, switched by an error function with parameter $\omega$ &
  $\omega$B97X, $\omega$B97M-V \\
double hybrid & the unoccupied orbitals, through a second-order
  perturbation (MP2) or random-phase correlation term & B2PLYP \\
\bottomrule
\end{tabular}
\end{table}

Dispersion, the weak attraction between non-bonded atoms, is missing from
all rungs and is added as a correction: D3~\cite{grimme2010} as an
empirical pairwise term, VV10~\cite{vydrov2010} as a functional of the
density. The suffix -D3 or -V in a functional name says which; a name
without suffix, such as $\omega$B97X, carries none.

Three levels of theory appear in this thesis:

\begin{itemize}
  \item \textbf{$\omega$B97X/6-31G(d)}~\cite{chai2008}: range-separated
    hybrid GGA without dispersion correction, small basis. The level of
    Transition1x, chosen for the cost of twelve thousand NEB runs and for
    compatibility with ANI-1x~\cite{schreiner2022}.
  \item \textbf{$\omega$B97M-V/def2-TZVP}~\cite{mardirossian2016}:
    range-separated hybrid meta-GGA with VV10 dispersion, triple-zeta
    basis. The target level of the delta correction. In the GMTKN55
    benchmark~\cite{goerigk2017} the functional belongs to the most
    accurate hybrids for thermochemistry, reaction barriers and
    non-covalent interactions.
  \item \textbf{$\omega$B97M-V/def2-TZVPD}: the same functional with
    diffuse functions added to the basis. The level of OMol25.
\end{itemize}


\subsection{Basis Sets}
\label{sec:background-basis}

The functional decides how the energy is computed from the density. The
basis set decides how well the orbitals, and with them the density, can be
represented at all.

In step 3 of the SCF iteration the orbitals have to be stored and
optimised. An orbital is a function of the electron coordinates, and a
computer cannot handle an arbitrary function. Each orbital is therefore
written as a linear combination of a fixed set of known functions
$\chi_\mu$ centred on the atoms,
%
\begin{equation}
  \phi_i(\mathbf{r}) = \sum_\mu c_{\mu i}\, \chi_\mu(\mathbf{r}) .
\end{equation}
%
The set $\{\chi_\mu\}$ is the basis set. Its functions are fixed before the
calculation; the SCF adjusts only the coefficients $c_{\mu i}$, that is,
how much of each basis function enters an orbital. An orbital can therefore
only take shapes that can be mixed from the basis functions. Gaussian
functions $e^{-\alpha r^2}$ are used because the integrals over them can be
solved analytically; a basis function is usually a fixed sum of several
Gaussians of different width~\cite{jensen2017}.

The size of the basis set is a trade-off. A small basis restricts the
orbitals to shapes close to those of the free atoms and is cheap. A large
basis lets the orbitals adapt to the molecule and approaches the basis set
limit, at a cost that grows formally as $\mathcal{O}(N^4)$ in the number of
functions.

A basis set is described by three properties:

\begin{itemize}
  \item \textbf{Functions per valence orbital.} Each valence orbital of
    an atom is built from one or more basis functions of different
    width. With one function the orbital keeps the size it has in the
    free atom. With two (double zeta) the SCF can mix a narrow and a
    wide function, so the orbital can shrink or expand in the molecule;
    with three (triple zeta) it has more freedom still. Core orbitals
    hardly change on bonding and usually get one function.
  \item \textbf{Polarisation functions.} Functions of higher angular
    momentum than the atom's occupied orbitals, $d$ on C, N, O and $p$ on
    H. They let an orbital bend away from the atom, which a bond requires.
  \item \textbf{Diffuse functions.} Functions that extend far from the
    nucleus. They are needed for anions and weakly bound electrons.
\end{itemize}

Three basis sets appear in this thesis (Table~\ref{tab:basis}). Each is
larger than the previous one in one of these properties.

\begin{table}[h]
\centering
\caption{Basis sets used in this thesis.}
\label{tab:basis}
\begin{tabular}{@{}lccc@{}}
\toprule
 & 6-31G(d)~\cite{hariharan1973} & def2-TZVP~\cite{weigend2005} & def2-TZVPD~\cite{rappoport2010} \\
\midrule
functions per valence orbital & 2 & 3 & 3 \\
polarisation functions        & on C, N, O & on all atoms & on all atoms \\
diffuse functions             & no & no & yes \\
functions per C atom          & 15 & 31 & 37 \\
\midrule
used for & Transition1x & delta target level & OMol25 \\
\bottomrule
\end{tabular}
\end{table}

Functional and basis set together are the level of theory, written
functional/basis, for example \mbox{$\omega$B97X/6-31G(d)}. The two
errors are independent: a larger basis does not correct the functional,
and a better functional does not correct a small basis. A third choice
remains to define a Kohn--Sham calculation: the treatment of spin.


%%%%%%%%%%%%%%%%%%%%%%%%%%%%%%%%%%%%%%%%%%%%%%

\subsection{Spin States of a Molecule}
\label{sec:background-spinstates}

Each electron has spin up ($\alpha$) or spin down ($\beta$). The spins of
all electrons add to the total spin $S$ of the molecule; two electrons of
opposite spin cancel. The input of a calculation specifies the
multiplicity $2S + 1$. All molecules in this thesis have $S = 0$,
multiplicity one, and are called singlets.

A singlet can have two different electronic structures:

\begin{itemize}
  \item \textbf{closed-shell:} every occupied orbital holds one $\alpha$
    and one $\beta$ electron. This is the normal case for a stable
    molecule.
  \item \textbf{open-shell singlet:} two electrons of opposite spin sit in
    two different orbitals, for example one on each end of a breaking
    bond. This is the case at diradicals and at many transition states.
\end{itemize}

Both are entered as multiplicity one. The input does not say which one is
meant. Whether a calculation can find the open-shell structure depends on
how it treats spin, which is the subject of the next section.

\subsection{Restricted and Unrestricted Kohn--Sham}
\label{sec:background-spin}

The presentation follows Szabo and Ostlund~\cite{szabo1996}.

\paragraph{Two treatments.} The third choice of Figure~\ref{fig:choices}
is whether an $\alpha$ electron must share its spatial orbital with a
$\beta$ electron (restricted) or may have its own (unrestricted). This
decides which of the two structures of
Section~\ref{sec:background-spinstates} the calculation can describe
(Table~\ref{tab:rks-uks}). Quantum chemistry programs such as ORCA and
Q-Chem run RKS for multiplicity one unless told
otherwise~\cite{orca_manual, qchem_manual}.

\begin{table}[h]
\centering
\small
\caption{The two spin formalisms and what they can describe.}
\label{tab:rks-uks}
\begin{tabular}{@{}lcc@{}}
\toprule
 & RKS & UKS \\
\midrule
orbitals & $\phi_i^{\alpha} = \phi_i^{\beta}$ & independent \\
energy & & $\le$ RKS \\
program default for multiplicity one & yes & \\
\midrule
closed-shell & $\langle S^2 \rangle = 0$ & $\langle S^2 \rangle = 0$ \\
open-shell singlet & not possible & $\langle S^2 \rangle > 0$, broken symmetry \\
\bottomrule
\end{tabular}
\end{table}

\paragraph{$\langle S^2 \rangle$: which structure was found.} A UKS
calculation can end in either structure. The number that tells which is
$\langle S^2 \rangle$, the expectation value of the total spin operator,
which every unrestricted calculation reports. Its exact value is $S(S+1)$,
so 0 for any singlet. A single determinant gives:

\begin{itemize}
  \item $\langle S^2 \rangle = 0$ when every $\alpha$ orbital equals its
    $\beta$ partner: the closed-shell structure;
  \item $\langle S^2 \rangle \approx 1$ when one $\alpha$ and one $\beta$
    electron sit in orbitals that do not overlap: the open-shell structure.
    The value is not 0 because a single determinant cannot represent an
    open-shell singlet exactly; it is an equal mixture of singlet and
    triplet;
  \item values in between for partial separation.
\end{itemize}

The excess over 0 is called spin contamination. It is a property of the
determinant, not of the molecule. For a singlet, $\langle S^2 \rangle > 0$
means the calculation has found the open-shell structure.

\paragraph{Broken symmetry.} A UKS singlet solution with
$\langle S^2 \rangle > 0$ is called a broken-symmetry (BS)
solution~\cite{noodleman1981}. The name says what happened: the symmetry
between $\alpha$ and $\beta$ orbitals is broken, while the total spin stays
zero. It is the single-determinant description of the open-shell singlet.
Whether such a solution exists depends on the geometry. For H$_2$ at its
equilibrium distance it does not; both electrons share the bonding
orbital. Beyond a certain bond length it appears, lies below the RKS
solution, and has one electron on each atom. The same happens at
transition states in which a bond is broken~\cite{bursch2022}.

\subsection{Breaking the Symmetry and Testing for It}
\label{sec:background-stability}

\paragraph{Symmetry does not break by itself.} Requesting UKS is not
enough to obtain a broken-symmetry solution. If the starting orbitals are
the same for both spins, both spins see the same potential at every SCF
step, and the iteration converges to the RKS solution even where a lower
one exists. The symmetry has to be broken on purpose. Two ways:

\begin{itemize}
  \item \textbf{Perturb the starting guess}, typically by mixing the
    $\beta$ HOMO and LUMO. The SCF then starts outside the symmetric
    solution and usually converges to the lower one. It can also fall back
    to the RKS solution; a closed-shell result then does not say whether a
    lower solution exists.
  \item \textbf{Test the converged solution} with a stability analysis. It
    decides whether a lower solution exists, and gives the direction to
    it.
\end{itemize}

\paragraph{Stability analysis~\cite{cizek1967, seeger1977}.} A converged
SCF solution is a stationary point of the energy with respect to rotations
of the orbitals. Whether it is a minimum is decided by the second
derivative of the energy with respect to these rotations:

\begin{enumerate}
  \item Converge the SCF.
  \item Compute the second derivative of the energy with respect to
    orbital rotations.
  \item Examine its lowest eigenvalue. A negative eigenvalue in the
    direction from restricted to unrestricted orbitals means a lower UKS
    solution exists.
  \item Follow that eigenvector and reconverge. This gives the
    broken-symmetry solution; the energy difference
    $E_\mathrm{RKS} - E_\mathrm{BS}$ is the breaking depth.
\end{enumerate}

\paragraph{Terminology.} A structure is \textbf{closed-shell} when the
stability analysis confirms the RKS solution as the lowest,
$\langle S^2 \rangle = 0$, and \textbf{broken-symmetry} when a lower UKS
solution exists, $\langle S^2 \rangle > 0$. In the first case the
restricted solution is called \textbf{stable}, in the second
\textbf{unstable}~\cite{seeger1977}.

\paragraph{Two surfaces.} Each spin formalism defines its own potential
energy surface. The \textbf{restricted surface} is the RKS energy at
every geometry. The \textbf{unrestricted surface} is the lowest UKS
solution at every geometry, found with the stability analysis: where
the restricted solution is stable the two surfaces coincide, where it
is unstable the unrestricted surface lies below. The stability analysis
is a statement about the electronic solution at one fixed geometry.
Whether that geometry is also a stationary point of the unrestricted
surface is a separate question, answered by the gradient of the
unrestricted solution at that geometry.


% %%%%%%%%%%%%%%%%%%%%%%%%%%%%%%%%%%%%%%%%%%%%%%
\subsection{Beyond Kohn--Sham DFT}
\label{sec:background-beyond}

The presentation follows Jensen~\cite{jensen2017}, Chapter 4. This work
considered a coupled cluster method as target level and attempted a
multireference method for reference geometries; neither enters the
final results. They are described here as far as
needed to follow why.

\paragraph{CCSD(T).} DFT approximates the correlation between electrons
with the functional. Coupled cluster calculates it from the wavefunction.
\mbox{CCSD(T)}~\cite{raghavachari1989}, coupled cluster with single,
double and perturbative triple excitations, is the standard variant. For
single-reference molecules the error is below 1~kcal/mol. It is the
standard reference for testing density functionals. The cost scales as
$\mathcal{O}(N^7)$ with the number of basis functions, for gradients as
well~\cite{scuseria1991}, so it is used for single-point energies and not
for path searches. It is a single-reference method and fails in the
multireference case.

\paragraph{CASSCF and NEVPT2.} A single determinant assigns every
electron to one orbital. In the multireference case the wavefunction
is a combination of several such assignments with comparable weights:
for a breaking bond, both electrons in the bonding orbital and both in
the antibonding one. No single determinant is then a good starting
point, and the combination has to be built in from the start. Only a
few orbitals change their occupation between the determinants that
mix, here the bonding and the antibonding orbital of that bond. These
are the active orbitals. Distributing the active electrons over them
in every possible way, with all lower orbitals doubly occupied and all
higher ones empty, gives the active space. CASSCF starts from all its
determinants at once,
%
\begin{equation}
  \Psi_\text{CAS} = \sum_{I \in \text{active}} c_I \,\Phi_I ,
\end{equation}
%
and optimises the weights $c_I$ and the orbitals together. This
captures the mixing of the few large determinants, but not the many
small contributions from the orbitals outside the active space.
NEVPT2~\cite{angeli2001} adds those afterwards, by perturbation theory
on top of the CASSCF wavefunction. Both results depend on which
orbitals are active. For a single geometry that is a
matter of choice; for a reaction path it is a condition: the same orbitals
must stay active at every geometry, otherwise the surface changes its
definition along the path. Choosing the active space at one geometry is
automated~\cite{sayfutyarova2017, stein2019autocas}; keeping it consistent
along a path is not~\cite{bensberg2023, seal2025wasp}.

\paragraph{Multireference diagnostics.} Whether a molecule is a
multireference case is not known in advance. Diagnostics estimate it from
a cheaper calculation. None of them measures the multireference character
directly; each measures a quantity that grows with it:

\begin{itemize}
  \item the $T_1$ diagnostic~\cite{lee1989}: the size of the single
    excitation amplitudes of a CCSD calculation. Values above about 0.02
    indicate that one determinant is not enough. It needs a CCSD
    calculation and is therefore as expensive as the method it is meant
    to replace.
  \item the $M$ diagnostic~\cite{tishchenko2008}: computed from the
    occupation numbers of a CASSCF wavefunction. It needs a CASSCF
    calculation and inherits its dependence on the active space.
  \item $N_\text{FOD}$~\cite{grimme2015}: computed from one DFT
    calculation in which the orbital occupations are smeared at an
    electronic temperature of 5000~K. Orbitals that are close in energy
    then take fractional occupations $n_i$, and
    %
    \begin{equation}
      N_\text{FOD} = \sum_i |n_i - n_i^0|
    \end{equation}
    %
    sums the deviations from the integer occupations $n_i^0$. Near-zero
    means one determinant suffices; larger values mark a multireference
    candidate.
\end{itemize}

All three characterise a structure, not a particular calculation: they
estimate whether one determinant is likely to be enough there.

% =====================================================================
\section{The Nudged Elastic Band Method}
\label{sec:background-neb}

An optimiser that follows the force downhill ends in a minimum. A
saddle point is a maximum along one direction, so a downhill search
moves away from it, and it has to be found another way. If reactant
and product are known, a path between them can be optimised instead: the
minimum energy path, the path of lowest energy connecting the two minima,
whose highest point is the transition state. The nudged elastic band
(NEB)~\cite{jonsson1998, henkelman2000tangent} represents the path as a
chain of images $\mathbf{R}_0, \dots, \mathbf{R}_{N_\text{img}-1}$ with the
endpoints fixed at reactant and product. Neighbouring images are connected
by springs. On each image the true force acts only perpendicular to the
path and the spring force only along it,
%
\begin{equation}
  \mathbf{F}_i = -\nabla E(\mathbf{R}_i)\big|_\perp
               + \mathbf{F}_i^\text{spring}\big|_\parallel ,
\end{equation}
%
so the images relax onto the minimum energy path and stay evenly spaced on
it.

The climbing image variant (CI-NEB)~\cite{henkelman2000cineb} treats the
highest image differently: its spring force is removed and the component
of the true force along the path is inverted,
%
\begin{equation}
  \mathbf{F}_i^\text{CI} = -\nabla E(\mathbf{R}_i)
  + 2\,\nabla E(\mathbf{R}_i)\big|_\parallel .
\end{equation}
%
This image moves uphill along the path and downhill perpendicular to
it, and converges to the saddle point. The two are run in sequence: a
plain NEB first, until the band lies roughly on the minimum energy
path, then the climbing image is switched on and the band is converged
to the final tolerance. Every search in this thesis follows this
two-stage protocol (Sections~\ref{sec:delta-protocol} and
\ref{sec:neb}).

The optimisation stops when the largest component of the projected force,
over all images and all $3N$ Cartesian directions of the $N$ atoms, falls
below a threshold $f_\text{max}$, typically 0.05~eV\,\AA$^{-1}$. The force
that remains at the end is the residual force.

A converged CI-NEB gives a transition state and a barrier on the surface
it was run on. The transition state is a maximum along the path and a
minimum perpendicular to it. That it has exactly one negative Hessian
eigenvalue is not verified unless a frequency calculation follows.


% =====================================================================
\section{Machine-Learned Interatomic Potentials}
\label{sec:background-mlip}

The presentation follows the reviews of Behler~\cite{behler2021} and Unke
et al.~\cite{unke2021}.

\subsection{The Approximation}
\label{sec:mlip-approximation}

A machine-learned interatomic potential (MLIP) replaces the electronic
structure calculation by a function fitted to its output. The energy is a
sum over atoms,
%
\begin{equation}
  E(\mathbf{R}) = \sum_a \varepsilon_a(\mathcal{N}_a),
  \label{eq:mlip-decomposition}
\end{equation}
%
where $\varepsilon_a$ is the contribution of atom $a$ and depends only on
its environment $\mathcal{N}_a$, the atoms within a cutoff radius of
typically 5 to 6~\AA. The forces are the analytic gradient of
Eq.~\eqref{eq:mlip-decomposition}. Two consequences of this form:

\begin{itemize}
  \item the cost grows linearly with the number of atoms, and a model
    trained on small molecules can be applied to larger ones;
  \item the model has no wavefunction, no density and no spin.
    It represents one function of the nuclear coordinates. Which surface
    that function approximates is fixed by the data it was fitted to.
\end{itemize}

\subsection{Symmetry}
\label{sec:mlip-symmetry}

The energy of a molecule does not change when the molecule is translated,
rotated, or when identical atoms are exchanged; the forces rotate with the
molecule. A model must respect this, otherwise it predicts different
energies for the same structure in different orientations. Two ways:

\begin{itemize}
  \item \textbf{Invariant models} describe the environment by distances
    and angles, which do not change under rotation, and any function of
    them is invariant automatically.
  \item \textbf{Equivariant models} carry features that transform under
    rotation in a defined way, scalars, vectors and higher tensors, and
    read the energy out from the scalars only
    (Section~\ref{sec:mlip-mace}). They reach a given accuracy with fewer
    training data, and all models in this thesis are of this type.
\end{itemize}

\subsection{Message Passing}
\label{sec:mlip-mace}

\paragraph{Features.} Every atom carries a list of numbers, its features,
that describes the atom itself and its surroundings. At the start the
list contains only the element of the atom. Over the layers described
next, it is filled with information about the neighbours: how many there
are, how far away, in which directions, and what their surroundings look
like in turn. The list is a compressed description of the local structure
around the atom, from which its energy contribution can be read off.
Which properties the list records is not fixed by hand; the network
adjusts this during training so that the energies and forces of the
training data are reproduced.

\paragraph{Message passing.} The list is filled in layers. In every
layer an atom receives from each neighbour within the cutoff the
neighbour's features together with the distance to it, a message, and
updates its own features from them. After one layer an atom knows its
direct neighbours, after two layers also their neighbours. Early models
passed only distances, so an atom learned how far its neighbours are
but not in which directions. PaiNN~\cite{schutt2021painn} added the
directions, stored as vectors that rotate with the molecule; such
features are called equivariant (next paragraph).
MACE~\cite{batatia2022} lets one message combine several neighbours at
once instead of one at a time, so that two layers are enough.

\paragraph{Channels.} The feature list of an atom is divided into
channels. Each channel records one learned property of the environment,
and it holds as many numbers as that kind of property needs
(Table~\ref{tab:channels}). The kinds are labelled $\ell$, the same label
that distinguishes $s$, $p$, $d$ and $f$ orbitals. The network has many
channels of each kind. The MACE model of this thesis has 1024 channels
for every $\ell$ from 0 to 3 (Table~\ref{tab:mace-settings}); a channel
of kind $\ell$ holds $2\ell + 1$ numbers, so the list of one atom holds
$1024 \times (1 + 3 + 5 + 7) = 16\,384$ numbers per layer.

\begin{table}[h]
\centering
\caption{Kinds of channels in an equivariant model. The examples are
illustrations; the network decides during training what each channel
records.}
\label{tab:channels}
\begin{tabular}{@{}clp{4.6cm}p{3.6cm}@{}}
\toprule
$\ell$ & numbers per channel & example of a property & under rotation \\
\midrule
0 & 1 & how crowded the surroundings are & unchanged \\
1 & 3 & mean direction to the neighbours & rotates like a vector \\
2 & 5 & whether the neighbours lie in a plane & rotates like a $d$ orbital \\
3 & 7 & a more detailed shape & rotates like an $f$ orbital \\
\midrule
\multicolumn{4}{@{}l}{1024 channels of each kind: $1024 \times 16 = 16\,384$ numbers per atom and layer} \\
\bottomrule
\end{tabular}
\end{table}

\paragraph{Readout.} A readout block is a small network that maps the
features of an atom to its energy contribution $\varepsilon_a$ in
Eq.~\eqref{eq:mlip-decomposition}. Because $\varepsilon_a$ is a scalar
and cannot depend on a direction, the readout uses only the $\ell = 0$
channels. The vector and tensor channels matter inside the layers, where
their combinations produce the scalar channels of the next layer.

\paragraph{Intuition.} An illustration, not a description of the
computation. A carbon atom in a C--C bond stretched to 2.0~\AA{} receives
from its neighbour a message built from the neighbour's element and
learned functions of the distance; at 2.0~\AA{} these produce a different
pattern across the channels than at 1.54~\AA. Only the atoms within the
cutoff enter, so a carbon with the same surroundings in a different
molecule gets the same pattern. The readout has learned that this pattern
goes with a higher energy. No channel is a named property; the channels
record whatever helped to reproduce the training energies.

\subsection{Conservative and Direct Forces}
\label{sec:mlip-conservative}

Forces can be obtained in two ways. A conservative model predicts the
energy and takes the forces as its gradient, as in
Eq.~\eqref{eq:mlip-decomposition}. A direct model predicts the forces with
a separate output. Direct forces are cheaper and often more accurate on
force benchmarks, but they are not the gradient of any energy: the work
along a closed path does not vanish, and a point of zero predicted force
is not a stationary point of the predicted energy. A saddle-point search
and a barrier height both require forces that are the gradient of the
reported energy. All models in this thesis are conservative.

\subsection{Training}
\label{sec:mlip-training}

The presentation follows Goodfellow et al.~\cite{goodfellow2016}.

\paragraph{Loss.} A model is fitted to reference energies and forces by
adjusting its weights to minimise
%
\begin{equation}
  L = \lambda_E \, d\big(E_\text{pred}, E_\text{ref}\big)
    + \lambda_F \, d\big(\mathbf{F}_\text{pred}, \mathbf{F}_\text{ref}\big),
\end{equation}
%
summed over the training structures. The function $d$ measures the
distance between prediction and reference,
$\Delta = x_\text{pred} - x_\text{ref}$. Three choices are common:
%
\begin{align}
  d_\text{sq}(\Delta)  &= \Delta^2, \\
  d_\text{abs}(\Delta) &= |\Delta|, \\
  d_\text{Huber}(\Delta) &=
  \begin{cases}
    \tfrac{1}{2}\Delta^2 & |\Delta| \le \delta \\
    \delta\,\big(|\Delta| - \tfrac{1}{2}\delta\big) & |\Delta| > \delta .
  \end{cases}
\end{align}
%
The squared loss weights large errors heavily, so a few structures with
large errors can dominate the sum. The absolute loss treats all errors
alike but has no smooth minimum. The Huber loss is quadratic below the
threshold $\delta$ and linear above it, and combines the two. The
weights $\lambda_E$ and $\lambda_F$ set how much energies and forces
count. Forces dominate the fit: each structure contributes $3N$ force
components and one energy, and the forces fix the local shape of the
surface, which a geometry optimisation follows.

\paragraph{Data splits.} The available structures are divided into three
sets (Table~\ref{tab:training-terms}). The training set is what the
weights are fitted to. The validation set is held out from the fit and
used during training to check whether the model still improves on
structures it has not seen, which decides when to lower the learning rate
and which checkpoint to keep. The test set is set aside before training
starts, is never used during it, and is evaluated once at the end to
report the accuracy on structures that played no role in training.

\paragraph{Procedure.} The loss is minimised by gradient descent: the
gradient of $L$ with respect to the weights is computed, and the weights
are moved a small step against it. Table~\ref{tab:training-terms} lists
the terms that describe this procedure.

\begin{table}[h]
\centering
\small
\caption{Terms of the training procedure.}
\label{tab:training-terms}
\begin{tabular}{@{}l p{9.4cm}@{}}
\toprule
term & meaning \\
\midrule
training set & the structures the weights are fitted to \\
validation set & structures held out from the fit; their loss is used to
  steer the training and to choose the final weights \\
test set & structures held out from both; used once, to report the
  accuracy \\
\midrule
batch & the structures whose loss is averaged for one weight update,
  typically 32 to 128 \\
epoch & one pass over the training set, or over a fixed random subsample
  of it if the set is too large \\
learning rate & the step size of the weight update \\
optimiser & the rule that turns the gradient into the step; Adam and
  AdamW~\cite{kingma2015adam, loshchilov2019adamw} adapt the step per
  weight from the history of its gradients \\
gradient clipping & an upper bound on the size of the gradient, to keep
  single batches from moving the weights too far \\
scheduler & a rule that lowers the learning rate during training; with
  ReduceLROnPlateau it is multiplied by a factor when the validation loss
  has not improved for a set number of epochs, the patience \\
checkpoint & the saved weights at one point of training; the final model
  is the checkpoint with the lowest validation loss \\
overfitting & the training loss keeps falling while the validation loss
  rises: the model fits the training structures rather than the surface \\
\bottomrule
\end{tabular}
\end{table}

\subsection{Raising the Level of Theory}
\label{sec:mlip-fidelity}

The accuracy of a model is bounded by the accuracy of its labels. A model
trained at one level of theory converges towards that level, not towards
the exact surface. Its errors against the exact surface are those of the
level plus those of the fit. Recomputing a full training set at a higher
level is usually not affordable. Three constructions avoid it:

\begin{itemize}
  \item \textbf{Delta learning}~\cite{ramakrishnan2015} keeps the
    low-level model and fits a second model to the difference between the
    two levels, $\Delta(\mathbf{R}) = E_\text{high}(\mathbf{R}) -
    E_\text{low}(\mathbf{R})$, at a sample of geometries. The prediction
    is the sum of both. The difference is smoother than either surface
    and needs far fewer training points.
  \item \textbf{Transfer learning}~\cite{smith2019ani1ccx} takes the
    low-level model as the starting point and continues training on
    high-level data, with all or part of the weights free.
  \item \textbf{Multifidelity learning} trains one model on both levels
    at once, with the level as an additional input.
\end{itemize}

\subsection{Limitations}
\label{sec:mlip-limits}

\begin{itemize}
  \item \textbf{Extrapolation.} An MLIP interpolates within the region of
    configuration space covered by its training data and has no defined
    behaviour outside it. Stretched bonds and transition states are
    sparsely covered by datasets of equilibrium structures.
  \item \textbf{No error estimate.} A model returns an energy and a force
    without a measure of confidence. A DFT calculation reports SCF
    convergence and can be tested for stability
    (Section~\ref{sec:background-stability}); an MLIP prediction carries
    no such signal.
\end{itemize}

% =====================================================================
\section{Datasets and Models}
\label{sec:background-datasets}

Two datasets and four models appear in this thesis. A dataset fixes two
things for every model trained on it: the level of theory the model
reproduces, and the region of configuration space it has seen. The
datasets are therefore described first.

\subsection{Dataset: Transition1x}
\label{sec:dataset-t1x}

Transition1x~\cite{schreiner2022} was built as training data for reactive
potentials: structures on and around reaction paths rather than at
equilibrium.

\paragraph{Reactions.} The starting point are the 11\,961 elementary
reactions of Grambow et al.~\cite{grambow2020}, each given as reactant,
transition state and product. The molecules contain up to seven heavy
atoms of C, N and O plus hydrogen, are neutral, closed-shell and in the
gas phase; the reactions are unimolecular rearrangements and
dissociations of these small organic molecules.

\paragraph{Geometries.} For every reaction a CI-NEB with ten images was
run at the \mbox{$\omega$B97X/6-31G(d)} level in ORCA. The level was
chosen for the cost of twelve thousand NEB runs and for compatibility
with ANI-1x. The band converged for
10\,073 reactions, with the criterion that the maximal perpendicular force
on the path does not exceed 0.05~eV\,\AA$^{-1}$. Every intermediate band
configuration of every NEB iteration was kept, not only the converged
path, so each reaction contributes on the order of 950 geometries. In
total the dataset holds 9.6~million structures with energies and forces.
Reactant, transition state and product of each reaction are stored
separately. All calculations are restricted.

\paragraph{Coverage.} The geometries lie on and near reaction paths of
one class of reactions at one level of theory. There is no solvent, no
charge or spin variation, and no element beyond C, H, N and O. Test and
validation data are 5\,\% of the total each.

\subsection{Dataset: OMol25}
\label{sec:dataset-omol}

Open Molecules 2025~\cite{omol25} is a dataset of more than 100 million
DFT calculations on about 83 million distinct molecular systems, built as
training data for general-purpose models.

\paragraph{Level of theory.} All calculations use
\mbox{$\omega$B97M-V/def2-TZVPD} in ORCA~6.0.0, with RI-J and COSX
integral approximations, tight SCF convergence and a large integration
grid. The dataset covers 83 elements, systems of 2 to 350 atoms with 50
on average, total charges from $-10$ to $+10$ and spin multiplicities
from 1 to 11. Non-singlets are computed unrestricted. Singlets are
computed unrestricted, with the initial guess perturbed by a $20^\circ$
rotation between HOMO and LUMO in the $\beta$ space, when they contain a
transition metal or lanthanide or when bonds are expected to break; all
other singlets are restricted. For the recomputed Transition1x structures
fewer than 5\,\% gave $\langle S^2 \rangle > 0.001$.

\paragraph{Composition.} The structures come from four chemical domains
and a fifth part that recomputes existing datasets
(Table~\ref{tab:omol25}). The geometries are taken from these sources as
they are, from molecular dynamics, optimisation trajectories, conformer
enumeration and reaction paths; no reaction path is relaxed with the
OMol25 protocol.

\begin{table}[h]
\centering
\small
\caption{Composition of OMol25~\cite{omol25}. Sizes are given in the
paper as atoms per domain.}
\label{tab:omol25}
\begin{tabular}{@{}l p{8.6cm} r@{}}
\toprule
domain & structures & atoms \\
\midrule
biomolecules & protein pockets and nucleic acids from the PDB; drug-like
  molecules from GEOM, ChEMBL and ZINC20 & 1.2\,B \\
electrolytes & molecules and ions of battery electrolytes in solution,
  water clusters of up to 70 molecules & 2.0\,B \\
metal complexes & complexes assembled from curated ligands and
  metal--oxidation-state pairs, in several spin states & 1.2\,B \\
main-group molecules & organic and heavy-main-group molecules and
  clusters, including reaction paths & -- \\
community & ANI-2x, ANI-1xBB, Transition1x, OrbNet Denali, SPICE2,
  Solvated Protein Fragments, about 30\,\% of GEOM, recomputed with the
  OMol25 protocol & 1.4\,B \\
\bottomrule
\end{tabular}
\end{table}

\paragraph{Reactions.} Transition-state regions enter through the
recomputed Transition1x, through metal-complex reactivity sets, through
electrolyte decomposition paths and through the interpolated and
artificial-force paths of the main-group domain. For Transition1x the
geometries are those of Section~\ref{sec:dataset-t1x}; only the energies
and forces are new.

\subsection{Model: MACE}
\label{sec:model-mace}

The MACE model of this thesis is trained on Transition1x. It has two
layers and 1024 channels per $\ell$ up to 3. The first layer keeps all
$\ell$; the second, last layer outputs scalars only. The features of both
layers are concatenated, giving 17\,408 numbers per atom, of which 2048
are scalars (Table~\ref{tab:mace-features}). A second readout block
attached to the same features, trained to a different target, is the
construction of Chapter~\ref{ch:delta}; the combined model is called
MACE+$\Delta$.

\begin{table}[h]
\centering
\caption{Per-atom features of the MACE encoder used in this thesis, 1024
channels per layer, $\ell \le 3$. Only the $\ell = 0$ column is
rotation-invariant.}
\label{tab:mace-features}
\begin{tabular}{@{}lccccr@{}}
\toprule
 & $\ell = 0$ & $\ell = 1$ & $\ell = 2$ & $\ell = 3$ & total \\
\midrule
layer 1 & 1024 & 3072 & 5120 & 7168 & 16\,384 \\
layer 2 & 1024 & --   & --   & --   & 1024 \\
\midrule
per atom & 2048 & 3072 & 5120 & 7168 & 17\,408 \\
\bottomrule
\end{tabular}
\end{table}

\subsection{Models: UMA and eSEN}
\label{sec:model-uma-esen}

Both models come from Meta FAIR and are trained on OMol25.

\paragraph{eSEN}~\cite{esen}, the equivariant Smooth Energy Network, is an
equivariant message-passing architecture built so that the energy is a
continuously differentiable function of the positions: no cutoff or
neighbour-list step introduces a jump. Its authors show that models
without this property, and models with direct forces, fail to conserve
energy in dynamics. The checkpoint used here is the small conserving
model trained on the full OMol25 dataset, eSEN-sm-conserving-all, with
6.3 million parameters and forces as the gradient of the
energy~\cite{omol25}.

\paragraph{UMA}~\cite{uma}, the Universal Model for Atoms, is built on
the eSEN architecture and trained on five datasets at once: OMol25,
OMat24 (materials), OC20 (catalysis), ODAC25 (direct air capture) and
OMC25 (molecular crystals), in total close to half a billion structures.
The five datasets were computed at different levels of theory, so the
model receives the dataset, called the task, as an input and reproduces
the level of that dataset; the task used in this thesis is \texttt{omol}.
Capacity is added by a mixture of linear experts: every linear layer
exists in several copies, and a weighted combination of them, chosen
from the composition, charge, spin and task of the structure, is used for
that structure. Only the combined layer runs during a calculation, so
the cost is that of a small model. Two sizes are used:

\begin{itemize}
  \item \textbf{UMA-S}: 150 million parameters in total, 6 million active
    per structure, 32 experts;
  \item \textbf{UMA-M}: 1.4 billion parameters in total, 50 million
    active.
\end{itemize}

Both predict forces as the gradient of the energy; the models are first
trained with direct forces and then fine-tuned to conservative ones. The
two sizes differ in capacity, not in training data or formulation.
UMA-S-1.2 was trained on the expanded OMol25 release, UMA-M-1.1 and the
eSEN checkpoint on the earlier one.

\subsection{Overview}
\label{sec:mlip-models}

\begin{table}[h]
\centering
\caption{The models of this thesis. All are equivariant message-passing
models with conservative forces.}
\label{tab:mlip-models}
\begin{tabular}{@{}llll@{}}
\toprule
model & training data & level of theory & role \\
\midrule
MACE & Transition1x & $\omega$B97X/6-31G(d) & base model, Chapter~\ref{ch:delta} \\
MACE+$\Delta$ & MACE plus correction head & $\omega$B97M-V/def2-TZVP & this work, Chapter~\ref{ch:delta} \\
UMA-S, UMA-M~\cite{uma} & OMol25 and others & $\omega$B97M-V/def2-TZVPD & Chapter~\ref{ch:mr} \\
eSEN~\cite{esen} & OMol25 & $\omega$B97M-V/def2-TZVPD & Chapter~\ref{ch:mr} \\
\bottomrule
\end{tabular}
\end{table}

MACE is trained on Transition1x alone and therefore specialised to reaction
paths of small organic molecules at one level of theory. UMA and eSEN are
general-purpose models trained on OMol25 and, for UMA, on further datasets
of materials, catalysts and molecular crystals.

% =====================================================================
\section{Software}
\label{sec:background-software}

Five packages are used in this thesis (Table~\ref{tab:software}). Their
division of labour follows a common pattern in computational chemistry:
one program computes energies and forces for a given geometry, another
decides which geometries to compute next.

\begin{table}[h]
\centering
\small
\caption{Software used in this thesis.}
\label{tab:software}
\begin{tabular}{@{}l l p{7.6cm}@{}}
\toprule
package & version & role \\
\midrule
ORCA~\cite{orca5} & 5.0.4 &
  DFT energies and gradients: single points, stability analysis,
  and the reference calculations behind the NEB paths \\
ASE~\cite{larsen2017ase} & 3.28.0 &
  Drives the geometry work: sets up the band, runs NEB and CI-NEB,
  relaxes endpoints, and calls ORCA or an MLIP for every energy and
  force it needs \\
PySCF~\cite{pyscf} & 2.12.1 &
  $N_\text{FOD}$ screening (DFT with Fermi smearing) and the
  multireference attempt (CASSCF, NEVPT2) \\
MACE~\cite{mace_code} & 0.3.15 &
  Training of the MACE base model and of the correction head; MLIP
  inference for MACE and MACE+$\Delta$ \\
fairchem~\cite{fairchem} & 2.20.0 &
  Inference for the UMA and eSEN models \\
\bottomrule
\end{tabular}
\end{table}

\paragraph{ORCA} is a quantum chemistry program for molecules, with a
focus on DFT and correlated wavefunction methods. In this thesis it
performs every DFT calculation: the single points on model geometries,
the stability analysis and the reconvergence to the broken-symmetry
solution, and the energies and gradients of the reference NEB paths.
Density fitting for the Coulomb and exchange integrals (RIJCOSX) keeps
hybrid functionals affordable at triple-zeta basis size. ORCA is also the
program with which OMol25 was computed, in a later version.

\paragraph{ASE}, the Atomic Simulation Environment, is a Python library
for setting up, running and analysing atomistic calculations. It holds
the geometries, provides the optimisers, BFGS for minima and the NEB
driver with the ode12r integrator~\cite{makri2019preconditioning} for
paths, and applies the convergence criterion. It does not compute
energies itself. Every source of energies and forces is attached as a
calculator: a geometry goes in, energy and forces come out, and ORCA,
MACE and the fairchem models all take this role. A NEB run therefore uses
the same driver, the same optimiser and the same convergence criterion
whether the forces come from DFT or from a model; only the calculator
differs. This is what makes the comparison between DFT-driven and
model-driven searches in this thesis a comparison of the surfaces, not
of the search procedure.

\paragraph{PySCF} is a Python-based quantum chemistry package in which
the individual steps of a calculation are accessible as functions. It is
used where a calculation departs from a standard protocol: the
$N_\text{FOD}$ screening, which needs DFT with Fermi smearing of the
occupations, and the multireference attempt, which needs CASSCF with an
automated active-space selection.

\paragraph{MACE and fairchem} are the two model packages. The MACE
package trains and evaluates MACE models and exposes the per-atom
features of a trained model, which the correction head of
Chapter~\ref{ch:delta} reads. fairchem is the library of Meta FAIR that
ships the UMA and eSEN checkpoints; a model is selected by checkpoint
name and, for UMA, by the task head, here the molecular task trained on
OMol25. Both packages provide an ASE calculator.
